% Mechanical relocation generated from the preservation ledger.
% Existing prose is included byte-for-byte from preserved blocks.

\section{Persistence, Restart, and Reporting}
\subsection{Persistent Artifacts and Recovery}

\paragraph{Persistent assignment artifacts.}
Assignment preparation is also demand-independent. Its immutable reusable state
consists of the time-expanded graph arrays and labels, destination/OD grouping
arrays, and generalized-cost components. It does not contain demand magnitudes,
observations, priors, optimizer settings, or compiled executables. The package
can store this state as compressed NPZ plus canonical JSON through optional
\texttt{off}, \texttt{auto}, \texttt{refresh}, and \texttt{readonly} policies;
ordinary uncached calls remain the default.

The assignment-cache provenance includes the complete graph-relevant scenario
contents, OD keys, assignment and cost configuration, numeric type, package and
schema versions, sentinel conventions, and indexing rules. Loaded archives are
checked for exact provenance, recognized fields, counts, shapes, dtypes, CSR and
padded-mask consistency, and integer/Boolean conventions before JAX arrays are
reconstructed. Cache writes use a unique temporary file, synchronization to
disk, and atomic replacement. Thus concurrent writers are safe, corrupt or
incompatible entries rebuild under \texttt{auto}, and \texttt{readonly} never
writes. The measurement-operator cache retains its own compact-input content
hash, so a valid assignment archive cannot validate an incompatible operator.

Internal TPG profiling initially assigned 3.142 s of a 3.313 s preparation to
destination grouping. The cause was repeated Python scalar extraction from JAX
node arrays, which introduced thousands of small device synchronizations. One
NumPy view is now reused for immutable node indexing and validation, and all
destination link masks are constructed in one broadcasted operation. Uncached
assignment construction consequently fell to 0.097 s: graph construction took
about 0.049 s, OD indexing and grouping 0.022 s, and cost construction plus
synchronization 0.026 s in the compilation-cache-enabled run.

The TPG artifact contains 36,478 nodes, 11,259 links, 22,022 OD records, and 113
destination groups. Its explicit arrays occupy 6,909,418 bytes and compress to
576,227 bytes. A representative strict cache hit spent 0.00044 s opening the
archive, 0.00460 s decompressing arrays, 0.01473 s in combined validation and
decompression, 0.000006 s reconstructing host objects, 0.00500 s transferring
and synchronizing JAX arrays, and 0.08787 s generating the canonical provenance.
The complete assignment-cache call took 0.1084 s.

\begin{table}[htbp]
\centering
\small
\begin{tabular}{lrrrrr}
\hline
Policy/case & Assignment & Preparation & Warm V\&G & Optimizer & Process \\
\hline
Off & 0.0971 & 1.1756 & 0.000992 & 0.2908 & 2.1939 \\
Populate & 0.3412 & 1.3532 & 0.000993 & 0.2975 & 2.4106 \\
Valid hit & 0.1084 & 1.0775 & 0.000937 & 0.2721 & 2.0211 \\
Corrupt/rebuild & 0.3201 & 1.3339 & 0.000965 & 0.3073 & 2.4081 \\
Mismatch/rebuild & 0.3321 & 1.3463 & 0.001048 & 0.3166 & 2.3935 \\
Read-only hit & 0.1128 & 1.1242 & 0.000972 & 0.2922 & 2.1872 \\
\hline
\end{tabular}
\caption{Fresh-process TPG assignment-cache measurements in seconds with the
measurement-operator and JAX compilation caches enabled.}
\label{tab:tpg-assignment-cache}
\end{table}

After eliminating scalar synchronization, uncached reconstruction is already
slightly faster than strict cache validation in the assignment stage itself;
canonical fingerprinting accounts for most of the hit time. The persistent
artifact is therefore useful primarily for strict reproducibility and validation,
while the focused grouping simplification provides the main speedup. Relative to
the original profile, complete one-iteration time falls from 5.389 s to roughly
2.0--2.2 s. Cache-hit runs at 1, 5, and 20 iterations reproduce objectives
55,243.4375, 49,635.4922, and 49,140.4258 exactly.


\paragraph{Anytime state and recovery.}
Initialization already constitutes a complete feasible approximate solution.
After every accepted block or conflict-free parallel batch, the estimator has a
complete flow vector, complete prediction, current and best objective, and
convergence diagnostics.  The durable journal records the replayable flow and
prediction deltas using temporary-file creation, synchronization, atomic
publication, and a separate commit marker.  Periodic compact checkpoints store
the complete state and exact next schedule position.  Resume is rejected when
any authoritative fingerprint changes, including scenario, assignment, OD
layout, fixed demand, measurements, prior, routing, partition, solver semantics,
or schema.  Returned statuses distinguish convergence, time, sweep, and update
budgets, graceful interruption with an approximate solution, resource guards,
and numerical failure.

Progress events do not claim a generic percentage precision.  They distinguish
exact, sampled, stale, deferred, and unavailable diagnostics and report objective
improvements, projected-gradient measures, block and variable coverage,
completed sweeps, elapsed time, checkpoint commitment, and estimated remaining
sweep time.

For a bounded large-network pilot, complete-operator products are governed by
an explicit fingerprinted policy.  The initial prediction may be computed or
supplied with validated shape, finiteness, bounds, fixed offset, and provenance.
Initial, periodic, and final exact projected gradients can be enabled
independently; a disabled diagnostic is recorded as deferred, never relabeled as
an exact global quantity.  Resume prediction validation may likewise be exact
or deferred.  Deferred operation avoids an otherwise mandatory global forward
or transpose product, but postpones the corresponding global consistency or
optimality claim.  A separate validation entry point can later recompute the
prediction, objective, and projected gradient, including in a fresh process.

A restricted pilot schedule is separately authorized from a complete sweep.
Every named block must belong to a completed support group, have compatible
persisted exact support, and satisfy support-row, nonzero, operator-storage, and
temporary-memory ceilings.  The checkpoint records that deterministic schedule
and its next position. Numerical semantics and the schedule remain
fingerprinted, whereas invocation limits may extend safely, so resuming a
five-update run with a six-update budget executes only the sixth update.

The monotonic elapsed-time deadline covers initialization, global products,
selected-block construction, solving, and checkpoint persistence.  JAX
compilation or dispatched device work cannot generally be interrupted inside a
process; an indivisible phase can therefore overshoot its allowance.  Such an
overshoot is reported explicitly, and no subsequent expensive phase is started.
Structured diagnostics expose global forward and transpose counts and times,
block construction and solve time, checkpoint time, prediction source, and
validation status.

The selected-block numerical kernels pass graph and routing arrays as explicit
dynamic JAX arguments rather than capturing complete assignment arrays in
Python closures.  Tracing, lowering, compilation, synchronized device
execution, and host transfer are timed separately.  A bounded lock-protected
compiled-executable cache is keyed by assignment provenance, backend, numeric
type, effective OD width, mapped-edge width, graph dimensions, and kernel schema;
deadline and callback state are excluded.  Deadline checks remain in Python
before and after each indivisible JAX stage.  On the public example,
fresh-process one-variable construction took 0.093--0.194 seconds, with zero
captured array constants and identical sparse results; compilation dominated
the roughly 0.001-second device execution.

For scheduler diagnostics, the builder can emit structured start and completion
events immediately around each phase.  An optional append-only JSON-lines sink
flushes every complete record and can call \texttt{fsync}; it is separate from
operator-cache validity.  Consequently, cancellation during XLA compilation
leaves a durable \texttt{jax\_compilation\_start} record without publishing an
incomplete numerical operator.  The public benchmark can stop intentionally
after tracing, lowering, compilation, or execution, allowing these indivisible
stages to be assigned separate external wall-time limits.



\subsection{Phase-Specific Reduced-OD Reuse}

Reduced-OD preprocessing artifacts are invalidated according to their actual
dependencies rather than one monolithic configuration fingerprint. Physical
stops, service-day selection, and timetable construction depend only on their
corresponding service inputs. Journey artifacts additionally depend on the
journey-search policy. Measurement responses and response operators add the
measurement policy and spatial output convention. Likelihood and production
mode belong to the model identity and therefore do not invalidate timetable,
journey, or response artifacts.

Every reusable artifact still validates its schema, phase fingerprint, parent
fingerprints, and numerical payload before acceptance. A changed likelihood
can consequently reuse the complete prepared response chain, whereas a changed
journey-search limit rebuilds journey choices and their descendants while
retaining compatible physical-stop, service, and timetable artifacts. A
changed production input republishes that input and the gravity/model stages
without rebuilding the measurement response operator. Corrupt or incompatible
artifacts fail closed.

Observed-data support reports are diagnostic artifacts, not scientific
operator identities.  Before a potentially long direct scheduled build, the
complete positive-boarding admission report is atomically written as
\texttt{positive\_boarding\_support\_preflight.json} in the identity-specific
checkpoint directory.  The report is refreshed at canonical-origin,
routing-support, and realized-cache stages.  Changing observations therefore
does not invalidate compatible routing or operator shards, but every new
activation rechecks the supplied positive boarding rows before reusing or
constructing an operator.

Potentially long journey construction accepts a progress callback that is not
part of any numerical fingerprint. Messages contain the phase and status,
completed and total query counts, current origin and departure, elapsed time,
recent mean query time, predicted remaining time, and peak resident memory.
Reporting is throttled by both a query interval and a wall-clock interval, with
the final query always emitted. The callback therefore provides useful live
feedback without changing reproducibility or flooding logs on fast cases.

The same reporting contract is used by long fixed-routing and temporal-cache
operations when their caller supplies a reporter.  Unit-based loops retain
only a bounded history of recent unit durations and use it to report
throughput, an estimated remaining time, a confidence level, and an expected
UTC completion time.  The first event may legitimately have no ETA; once a
measured unit has completed, the estimate is reported whenever the operation
exposes a meaningful unit of work.  Indivisible compilation or device calls
remain heartbeat-only when no defensible unit estimate is available.  Progress
events are throttled and additive: they do not trigger an extra objective,
gradient, routing, or JAX evaluation, and progress settings are excluded from
artifact and checkpoint identities.  A caller may explicitly disable a
reporter, but the high-level case-study workflow normally creates a durable
JSONL sink so that a long run does not appear silent.

Fixed-routing preparation distinguishes two independent reporting controls:
\path{progress_interval_seconds} limits the wall-clock spacing between
heartbeats, while \path{progress_interval_groups} requests updates after
completed destination groups or routing units.  Supplying both honors both
constraints; neither is converted into the other.  A terminal event is forced
after a completed phase, and a successful stage is accepted only after its
durable completion manifest and identity-specific artifact manifest validate.

Gravity run manifests use a separate model-level identity. They now contain
the complete serialized gravity specification and fingerprint, named parameter
blocks and transformations, feature-cache and direct-operator artifact
fingerprints, calibration, holdout, and unsupported-row summaries,
regularization settings, time-bin definitions, and destination-attractiveness
provenance. Checkpoint fingerprints include the specification and parameter
layout, so a changed scope, mapping, constraint, or ridge setting cannot resume
an incompatible optimizer state. Conversely, such a model change does not
invalidate a scientifically compatible fixed-routing assignment artifact.
